% problem_id: S001-lemma-etale-pullback-principal-divisor
% source_id: S001 (Stacks Project)
% source_file: spaces-chow.tex
% statement_locator: spaces-chow.tex lines 1462--1473
% proof_locator: spaces-chow.tex lines 1475--1586
% env: lemma
% id_note: label 在本来源内唯一
% pairing: tight (verified)
% status: complete (statement + author proof verbatim + reason + locators)
%
\begin{lemma}
\label{lemma-etale-pullback-principal-divisor}
In Situation \ref{situation-setup} let $f : X \to Y$ be an \'etale
morphism of good algebraic spaces over $B$. Assume $Y$ is integral.
Let $g \in R(Y)^*$. As cycles on $X$ we have
$$
f^*(\text{div}_Y(g)) =
\sum\nolimits_{X'} (X' \to X)_*\text{div}_{X'}(g \circ f|_{X'})
$$
where the sum is over the irreducible components of $X$
(Remark \ref{remark-irreducible-component}).
\end{lemma}

\begin{proof}
The map $|X| \to |Y|$ is open. The set of irreducible components of $|X|$
is locally finite in $|X|$. We conclude that $f|_{X'} : X' \to Y$
is dominant for every irreducible component $X' \subset X$.
Thus $g \circ f|_{X'}$ is defined (Morphisms of Spaces, Section
\ref{spaces-morphisms-section-rational-maps}), hence
$\text{div}_{X'}(g \circ f|_{X'})$ is defined. Moreover, the
sum is locally finite and we find that the right hand side indeed
is a cycle on $X$. The left hand side is defined by
Definition \ref{definition-flat-pullback}
and the fact that an \'etale morphism is flat of relative dimension $0$.

\medskip\noindent
Since $f$ is \'etale we see that $\delta_X(x) = \delta_y(f(x))$
for all $x \in |X|$. Thus if $\dim_\delta(Y) = n$, then $\dim_\delta(X') = n$
for every irreducible component $X'$ of $X$ (since generic points of $X$
map to the generic point of $Y$, see above). Thus both left
and right hand side are $(n - 1)$-cycles.

\medskip\noindent
Let $Z \subset X$ be an integral closed subspace with $\dim_\delta(Z) = n - 1$.
To prove the equality, we need to show that the coefficients
of $Z$ are the same. Let $Z' \subset Y$ be the integral closed
subspace constructed in Lemma \ref{lemma-proper-image}.
Then $\dim_\delta(Z') = n - 1$ too. Let $\xi \in |Z|$ be the generic point.
Then $\xi' = f(\xi) \in |Z'|$ is the generic point.
Consider the commutative diagram
$$
\xymatrix{
\Spec(\mathcal{O}_{X, \xi}^h) \ar[r] \ar[d] & X \ar[d] \\
\Spec(\mathcal{O}_{Y, \xi'}^h) \ar[r] & Y
}
$$
of Decent Spaces, Remark
\ref{decent-spaces-remark-functoriality-henselian-local-ring}.
We have to be slightly careful as the reduced Noetherian local rings
$\mathcal{O}_{X, \xi}^h$ and $\mathcal{O}_{Y, \xi'}^h$ need not be domains.
Thus we work with total rings of fractions $Q(-)$ rather than fraction fields.
By definition, to get the coefficient of $Z'$ in $\text{div}_Y(g)$
we write the image of $g$ in $Q(\mathcal{O}_{Y, \xi'}^h)$
as $a/b$ with $a, b \in \mathcal{O}_{Y, \xi'}^h$ nonzerodivisors
and we take
$$
\text{ord}_{Z'}(g) =
\text{length}_{\mathcal{O}_{Y, \xi'}^h}
(\mathcal{O}_{Y, \xi'}^h/a \mathcal{O}_{Y, \xi'}^h) -
\text{length}_{\mathcal{O}_{Y, \xi'}^h}
(\mathcal{O}_{Y, \xi'}^h/b \mathcal{O}_{Y, \xi'}^h)
$$
Observe that the coefficient of $Z$ in $f^*\text{div}_Y(G)$
is the same integer, see Lemma \ref{lemma-etale-pullback}.
Suppose that $\xi \in X'$. Then we can consider the maps
$$
\mathcal{O}_{Y, \xi'}^h \to
\mathcal{O}_{X, \xi}^h \to
\mathcal{O}_{X', \xi}^h
$$
The first arrow is flat and the second arrow is a surjective
map of reduced local Noetherian rings of dimension $1$.
Therefore both these maps send nonzerodivisors to nonzerodivisors
and we conclude the coefficient of $Z'$ in $\text{div}_{X'}(g \circ f|_{X'})$
is
$$
\text{ord}_Z(g \circ f|_{X'}) =
\text{length}_{\mathcal{O}_{X', \xi}^h}
(\mathcal{O}_{X', \xi}^h/a \mathcal{O}_{X', \xi}^h) -
\text{length}_{\mathcal{O}_{Y, \xi'}^h}
(\mathcal{O}_{X', \xi}^h/b \mathcal{O}_{X', \xi}^h)
$$
by the same prescription as above. Thus it suffices to show
$$
\text{length}_{\mathcal{O}_{Y, \xi'}^h}
(\mathcal{O}_{Y, \xi'}^h/a \mathcal{O}_{Y, \xi'}^h) =
\sum\nolimits_{\xi \in |X'|}
\text{length}_{\mathcal{O}_{X', \xi}^h}
(\mathcal{O}_{X', \xi}^h/a \mathcal{O}_{X', \xi}^h)
$$
First, since the ring map $\mathcal{O}_{Y, \xi'}^h \to
\mathcal{O}_{X, \xi}^h$ is flat and unramified, we have
$$
\text{length}_{\mathcal{O}_{Y, \xi'}^h}
(\mathcal{O}_{Y, \xi'}^h/a \mathcal{O}_{Y, \xi'}^h) =
\text{length}_{\mathcal{O}_{X, \xi}^h}
(\mathcal{O}_{X, \xi}^h/a \mathcal{O}_{X, \xi}^h)
$$
by Algebra, Lemma \ref{algebra-lemma-pullback-module}.
Let $\mathfrak q_1, \ldots, \mathfrak q_t$ be the nonmaximal
primes of $\mathcal{O}_{X, \xi}^h$ and set
$R_j = \mathcal{O}_{X, \xi}^h/\mathfrak q_j$.
For $X'$ as above, denote $J(X') \subset \{1, \ldots, t\}$
the set of indices such that $\mathfrak q_j$ corresponds
to a point of $X'$, i.e., such that under the surjection
$\mathcal{O}_{X, \xi}^h \to \mathcal{O}_{X', \xi}$
the prime $\mathfrak q_j$ corresponds to a prime
of $\mathcal{O}_{X', \xi}$.
By Chow Homology, Lemma \ref{chow-lemma-additivity-divisors-restricted}
we get
$$
\text{length}_{\mathcal{O}_{X, \xi}^h}
(\mathcal{O}_{X, \xi}^h/a \mathcal{O}_{X, \xi}^h) =
\sum\nolimits_j \text{length}_{R_j}(R_j/a R_j)
$$
and
$$
\text{length}_{\mathcal{O}_{X', \xi}^h}
(\mathcal{O}_{X', \xi}^h/a \mathcal{O}_{X', \xi}^h) =
\sum\nolimits_{j \in J(X')} \text{length}_{R_j}(R_j/a R_j)
$$
Thus the result of the lemma holds because
$\{1, \ldots, t\}$ is the disjoint union of the sets $J(X')$:
each point of codimension $0$ on $X$ lies on a unique $X'$.
\end{proof}
